\documentclass[10pt,a4paper]{amsart}
\usepackage[margin=19mm]{geometry}
\usepackage{amsmath,amssymb,mathtools}
\usepackage[scr=rsfs,cal=euler]{mathalfa}
\usepackage{xfrac}
\usepackage[unicode,hidelinks,bookmarks=false]{hyperref}
\newcommand{\ie}{i.e.\ }
\newcommand{\cf}{cf.\ }
\newcommand{\resp}{respectively\ }
\newcommand{\Subsection}{\subsection}
\providecommand{\nobreakdash}{\nobreak\hbox{-}}
\setlength{\emergencystretch}{2em}
\multlinegap0pt
\usepackage{bm}
\usepackage{bbm}
\usepackage{dsfont}
\usepackage{xspace}
\usepackage{cancel}
\usepackage{mathtools}
\usepackage[shortlabels]{enumitem}
\usepackage{amsthm}
\makeatletter
\@ifundefined{lem}{\newtheorem{lem}{Lemma}}{}
\makeatother
\title[Multi-to -one dimensional and semi-discrete screening]{Multi-to -one dimensional and semi-discrete screening}
\author{Omar Abdul Halim, Brendan Pass}
\date{Source: arXiv:2506.21740v2 (CC BY 4.0)}
\begin{document}
\maketitle

\noindent\emph{Source and licence.} Multi-to -one dimensional and semi-discrete screening. Version arXiv:2506.21740v2 (CC BY 4.0), distributed under the Creative Commons Attribution 4.0 International licence, as recorded in the arXiv licence metadata for this exact version and shown on the abstract page https://arxiv.org/abs/2506.21740v2. Full rights notice: licenses/arxiv-2506.21740-CC-BY-4.0.txt.\par\smallskip

\noindent\textit{Editorial scope.} Counted unit: the Lem whose source label is \texttt{case1} in the article (source0.tex lines 1075--1077, source label \texttt{case1}), delivered together with the article's own complete original proof of it (source0.tex lines 1078--1187, one \texttt{proof} environment of 110 lines). No mathematics is written, repaired or re-derived: statement and proof are the article's own text line for line. Required context retained uncounted: condition3 (lines 205--218). Every definition, display and label of those lines is retained unchanged. Omitted from this package and not redistributed: 1--204, 219--1074, 1188--1843. Nothing inside a retained excerpt is omitted or altered: every retained body is delivered exactly as the article's lines are written, so no omission receipt of any kind accompanies this package. Citations: the retained excerpts carry no citation command at all, so no bibliography entry of the article is transcribed and no reference is redistributed. Reference adaptation: none was needed and none was made; every reference inside the delivered unit resolves inside it, so no unresolved reference or citation is produced. Printed slips kept literal: no printed word, symbol, hypothesis or number of the counted statement or of its proof was altered. Layout and class: the article's own class is \texttt{amsart} with packages including graphicx, amsmath, amssymb, enumitem, geometry, comment, setspace, mathrsfs, amsthm, appendix, tikz, subcaption; the wrapper is the lane's pinned \texttt{amsart} editorial wrapper at 10pt on A4, which restates the theorem-like environments the retained text needs and adds the article's own macro definitions that the retained text needs (none), with extra wrapper packages (bm, bbm, dsfont, xspace, cancel, mathtools, enumitem). The delivered numbering is the unit's own. Assets: no retained span contains a figure environment, a graphics-inclusion command, a PDF-page inclusion, a table or a TikZ picture; no image file is copied into this package, the package declares no asset dependency and no image file is redistributed with it. Licence: version arXiv:2506.21740v2 of this article is distributed under the Creative Commons Attribution 4.0 International licence, as recorded in the arXiv licence metadata for this exact version and shown on the abstract page https://arxiv.org/abs/2506.21740v2. A full rights notice is delivered at licenses/arxiv-2506.21740-CC-BY-4.0.txt. Characters: every retained excerpt is pure ASCII, so the delivered engine, XeLaTeX with its default Latin Modern text font, reports no missing character.
\begin{enumerate}[label={(H\arabic*)}]
%\item\label{condition1} $\frac{c(z_j)-c(z_i)}{y_j-y_i} > \frac{c(z_i)-c(z_k)}{y_i-y_k},$
\item \label{condition2} $ \frac{c(z_j)-c(z_i)}{F(y_j)-F(y_i)}>\frac{c(z_i)-c(z_k)}{F(y_i)-F(y_k)},$\vspace{3pt}
  \item \label{condition3} $\frac{\frac{c(z_{i+1})-c(z_i)}{y_{i+1}-y_i}-\frac{c(z_i)-c(z_
  {i-1})}{y_i-y_{i-1}}}{\frac{F(y_{i+1})-F(y_i)}{y_{i+1}-y_i}-\frac{F(y_i)-F(y_{i-1})}{y_i-y_{i-1}}}>\frac{\frac{3}{2}\|f\|_\infty+\frac{\|f_{x_1}\|_\infty}{2}\Big(1+\frac{F(y_{i+1})-F(y_i)}{y_{i+1}-y_i} +\frac{c(z_{i+1})-c(z_i)}{y_{i+1}-y_i}\Big)}{\alpha},$
  
   \item \label{condition4} 
   \[\begin{array}{ll}
   \frac{F(y_{i+1})-F(y_{i})}{y_{i+1}-y_{i}}<&  \frac{2\alpha}{\|f\|_\infty\Big(2+\frac{\frac{F(y_{i+1})-F(y_{i})}{y_{i+1}-y_{i}}}{\frac{F(y_i)-F(y_{i-1})}{y_i-y_{i-1}}}\Big)+\|f_{x_2}\|_\infty\Big(1+2\frac{y_{i}-y_{i-1}}{F(y_{i})-F(y_{i-1})}+\frac{c(z_{i+1})-c(z_i)}{F(y_{i+1})-F(y_i)}\Big)}\times\\
   &\Bigg(1+\frac{\frac{c(z_{i+1})-c(z_i)}{F(y_{i+1})-F(y_i)}-\frac{c(z_{i})-c(z_{i-1})}{F(y_i)-F(y_{i-1})}}{\frac{y_{i}-y_{i-1}}{F(y_{i})-F(y_{i-1})}-\frac{y_{i+1}-y_i}{F(y_{i+1})-F(y_{i})}}\Bigg).
   \end{array}\]


\end{enumerate}

\begin{lem}\label{case1}
         Suppose that $u \in \mathcal{U}$ and $X_i =(Du)^{-1}(z_i)$ shares a boundary with only two other regions, $X_{i-1}$ and $X_{i+1}$.  Suppose that the two indifference curves intersect within $\overline{X}$, and that both intersect $[0,1]\times\{0\}$.  Then, under hypothesis   \ref{condition3}, $u$ is not a solution to the monopolist's problem.   
    \end{lem}

\begin{proof}We will show that lowering the prices of good $y_i$ by $\epsilon$ strictly increases profits.  Let $v_j = \max_{x \in X}x \cdot z_j - u(x)$ be the pricing schedule corresponding to $u$ and change $v_i \rightarrow v_i -\epsilon$, while leaving the other prices $v_j, j \neq i$ unchanged.

Letting $X_i^\epsilon$ be the region of consumers who choose good $z_i$ under the new price schedule.  We then have
\[\begin{array}{ll}
    \mathcal P(v^\epsilon) = \mathcal P(v)&-\int_{X_i^\epsilon} \epsilon f(x) dx\\
    &+ \int_{X_{i}^\epsilon \cap X_{i+1}}[x \cdot (z_{i}-z_{i+1}) - (u_{i}+\epsilon-u_{i+1}) -c(z_{i}) +c(z_{i+1}) ]f(x)dx \\
    &+ \int_{X_{i}^\epsilon \cap X_{i-1}} [x \cdot (z_{i}-z_{i-1}) - (u_{i}+\epsilon-u_{i-1}) -c(z_{i}) +c(z_{i-1})]
    f(x) dx 
\end{array}\]
where $u_j=u(x)$ for all $x\in X_j.$
 We differentiate this expression with respect to $\epsilon$.  Note that as $\epsilon$ varies, the region $X_i^\epsilon$ expands outward along its boundary curves 
$$
 L_i^\epsilon=X^N_=(y_i, v_{i+1} +\epsilon -v_{i})\cap \overline{ X^\epsilon_i} =\{x: x\cdot (z_{i+1}-z_i) =v_{i+1} +\epsilon -v_{i}\} \cap \overline{ X^\epsilon_i} \subseteq \overline{  X_{i}^\epsilon} \cap \overline{X_{i+1}}
 %\{x: x\cdot (z_{i+1}-z_i) =v_{i} -\epsilon -v_{i+1}\}
 $$ 
 and 
 $$
 L_{i-1}^\epsilon= X^N_=(y_{i-1}, v_{i} -\epsilon -v_{i-1})\cap \overline{ X^\epsilon_i} =\{x: x\cdot (z_{i}-z_{i-1}) =v_{i} -\epsilon -v_{i-1}\} \cap \overline{ X^\epsilon_i} \subseteq\overline{  X_{i}^\epsilon} \cap \overline{X_{i-1}}
 $$
 with outward unit normal speeds $\frac{1}{|z_i-z_{i+1}|}$ and $\frac{1}{|z_i-z_{i-1}|}$, respectively.  A standard formula from the calculus of moving boundaries then yields 
 \[\begin{array}{ll}
     &\frac{d}{d\epsilon} \mathcal P(v^\epsilon)  =\vspace{3pt}\\
     &-\int_{X_i^\epsilon}f(x)dx -\int_{L^\epsilon_i }\epsilon f(x)\frac{1}{|z_i-z_{i+1}|}d\mathcal H^{m-1}(x) 
     -\int_{L^\epsilon_{i-1}}\epsilon f(x)\frac{1}{|z_i-z_{i-1}|}d\mathcal H^{m-1}(x) \vspace{3pt}\\
     &-\int_{X_{i}^\epsilon \cap X_{i+1}}f(x)dx -\int_{X_{i}^\epsilon \cap X_{i-1}}f(x)dx\vspace{3pt}\\
     &+\int_{L^\epsilon_{i}}[x \cdot (z_{i}-z_{i+1}) - (u_{i}+\epsilon-u_{i+1}) -c(z_{i}) +c(z_{i+1}) ]f(x)\frac{1}{|z_i-z_{i+1}|}d\mathcal H^{m-1}(x)\vspace{3pt}\\
     &+
     \int_{L^\epsilon_{i-1}}[x \cdot (z_{i}-z_{i-1}) - (u_{i}+\epsilon-u_{i-1}) -c(z_{i}) -+(z_{i-1}) ]f(x)\frac{1}{|z_i-z_{i-1}|}d\mathcal H^{m-1}(x)
 \end{array}
 \]
Noting that $u_i=u_{i+1}$ and $u_i =u_{i-1}$ along the appropriate respective indifference curve, and that the volumes of the regions $X_{i}^\epsilon \cap X_{i+1}$ and $X_{i}^\epsilon \cap X_{i-1}$  dwindle to $0$ as $\epsilon \rightarrow 0$, we set $\epsilon =0$ to obtain
\begin{eqnarray}\label{optcond1}
         \frac{d}{d\epsilon}\Big|_{\epsilon=0} \mathcal P(v^\epsilon)  =
         &-&\int_{X_i}f(x)dx\label{derivative} \\
      &+&\int_{L^0_{i}}[x \cdot (z_{i}-z_{i+1})-c(z_{i}) +c(z_{i+1}) ]f(x)\frac{1}{|z_i-z_{i+1}|}d\mathcal H^{m-1}(x) \nonumber\\
&+&     \int_{L^0_{i-1}}[x \cdot (z_{i}-z_{i-1})-c(z_i) + c(z_{i-1}) ]f(x)\frac{1}{|z_i-z_{i-1}|}d\mathcal H^{m-1}(x) \nonumber
\end{eqnarray}
Now, let $a=(a_1,a_2)$ be the intersection point of the indifference regions $X^N_=(y_{i-1}, v_{i} - v_{i-1})$ and $X^N_=(y_i, v_{i+1} - v_{i})$.  





Since both indifference curves reach axis $ [0,1]\times\{0\} $ by assumption, we can parametrize them  $(x^{i-1}_1(x_2),x_2)$ and $(x^i_1(x_2),x_2)$ by $x_ 2 \in [0,a_2]$ and since the line segment $X^N_=(y_i, v_{i+1} - v_{i})$ is orthogonal to $z_{i+1}-z_i =( y_{i+1},F(y_{i+1})) -(y_{i},F(y_{i})) $, the slope of $x_1^i$ is $\frac{F(y_{i+1})-F(y_i)}{y_{i+1} -y_i}$, and $1$-dimensional Hausdorff measure (ie, arclength) along it is given by $d\mathcal H^{m-1}(x) = \sqrt{\big(\frac{F(y_{i+1})-F(y_i)}{y_{i+1} -y_i}\big)^2+1}dx_2$, so that 
\[\hspace{-0.4cm}\begin{array}{ll}
&\int_{L_i^0}[x \cdot (z_{i}-z_{i+1})-c(z_{i}) +c(z_{i+1}) ]f(x)\frac{1}{|z_i-z_{i+1}|}d\mathcal H^{m-1}(x)\vspace{3pt}\\
&=[a \cdot (z_{i}-z_{i+1})-c(z_{i}) +c(z_{i+1}) ]\frac{1}{|z_i-z_{i+1}|}\sqrt{\big(\frac{F(y_{i+1})-F(y_i)}{y_{i+1} -y_i}\big)^2+1}\int_0^{a_2} f(x^i_1(x_2),x_2)dx_2\vspace{3pt}\\
&=[a \cdot (z_{i}-z_{i+1})-c(z_{i}) +c(z_{i+1}) ]\frac{1}{y_{i+1}-y_{i}}\int_0^{a_2} f(x^i_1(x_2),x_2)dx_2\vspace{3pt}\\
&= [-a_1 -a_2 \frac{F(y_{i+1})-F(y_i)}{y_{i+1}-y_{i}} +\frac{c(z_{i+1}) -c(z_{i})}{y_{i+1}-y_{i}}]\int_0^{a_2} f(x^i_1(x_2),x_2)dx_2,
\end{array}\]
where  $a\cdot(z_{i+1}-z_{i})=x\cdot(z_{i+1}-z_{i})$ along $L_i^0.$
Similarly,
\begin{eqnarray*}
\int_{L^0_{i-1}}[x \cdot (z_{i}-z_{i-1})-c(z_i) + c(z_{i-1}) ]f(x)\frac{1}{|z_i-z_{i-1}|}d\mathcal H^{m-1}(x)\\
=[a_1 +a_2 \frac{F(y_{i})-F(y_{i-1})}{y_{i}-y_{i-1}} -\frac{c(z_{i}) -c(z_{i-1})}{y_{i}-y_{i-1}}]\int_0^{a_2} f(x_1^{i-1}(x_2),x_2)dx_2.
\end{eqnarray*}
%{\color{blue} BP: I added the lines in red below, which I think makes the argument more readable, and also corrected a handful of typos in the proof.  Please check you agree with it, and, if necessary, make similar changes in the proofs of the following Lemmas.}


 We can therefore rewrite \eqref{derivative} as 
\begin{equation}\label{eqn: derivative computation}
\begin{array}{ll}
         &\frac{d}{d\epsilon}\Big|_{\epsilon=0} \mathcal P(v^\epsilon)  =\vspace{3pt}\\
         &-\int_{X_i}f(x)dx -a_1\Big[\int_0^{a_2} f(x^i_1(x_2),x_2)dx_2-\int_0^{a_2} f(x_1^{i-1}(x_2),x_2)dx_2\Big] \vspace{3pt}\\
&-a_2\Big[\frac{F(y_{i+1})-F(y_i)}{y_{i+1}-y_{i}} \int_0^{a_2} f(x^i_1(x_2),x_2)dx_2 - \frac{F(y_{i})-F(y_{i-1})}{y_{i}-y_{i-1}}\int_0^{a_2} f(x_1^{i-1}(x_2),x_2)dx_2\Big]      \vspace{3pt}\\
&+     \frac{c(z_{i+1}) -c(z_{i})}{y_{i+1}-y_{i}}\int_0^{a_2} f(x^i_1(x_2),x_2)dx_2 -\frac{c(z_{i}) -c(z_{i-1})}{y_{i}-y_{i-1}}\int_0^{a_2} f(x_1^{i-1}(x_2),x_2)dx_2. 
\end{array}
\end{equation}



We bound the first term of \eqref{eqn: derivative computation} below by 
\begin{equation}\label{X_i mass bound}
-\int_{X_i}f(x)dx\ge -\|f\|_\infty\frac{a_2^2}{2}\Big(\frac{F(y_{i+1})-F(y_i)}{y_{i+1}-y_i}-\frac{F(y_i)-F(y_{i-1})}{y_i-y_{i-1}}\Big).
\end{equation}

Turning to the second term in \eqref{eqn: derivative computation}, since $a_1\le1,$ we have

\[
        \begin{array}{ll}
        &-a_1\Big(\int_0^{a_2}f(x_1^{i}(x_2),x_2)-f(x_1^{i-1}(x_2),x_2)dx_2\Big)\vspace{3pt}\\
        &\ge - \|f_{x_1}\|_\infty\int_0^{a_2}\Big(\frac{F(y_{i+1})-F(y_i)}{y_{i+1}-y_i}-\frac{F(y_{i})-F(y_{i-1})}{y_{i}-y_{i-1}}\Big)(a_2-x_2)dx_2\vspace{3pt}\\
        &\ge-\|f_{x_1}\|_\infty\Big(\frac{F(y_{i+1})-F(y_i)}{y_{i+1}-y_i}-\frac{F(y_{i})-F(y_{i-1})}{y_{i}-y_{i-1}}\Big)\frac{a_2^2}{2}.
        \end{array}\]
For the third term in \eqref{eqn: derivative computation}, we have 
\[\begin{array}{ll}
          &-a_2\Big(\frac{F(y_{i+1})-F(y_i)}{y_{i+1}-y_i}\int_0^{a_2}f(x_1^{i}(x_2),x_2)dx_2-\frac{F(y_i)-F(y_{i-1})}{y_i-y_{i-1}}\int_0^{a_2}f(x_1^{i-1}(x_2),x_2)dx_2\Big)\vspace{3pt}\\
        =&-a_2\Bigg(\frac{F(y_{i+1})-F(y_i)}{y_{i+1}-y_i}\int_0^{a_2}(f(x_1^{i}(x_2),x_2)-f(x_1^{i-1}(x_2),x_2))dx_2\vspace{3pt}\\
        &+\Big(\frac{F(y_{i+1})-F(y_i)}{y_{i+1}-y_i}-\frac{F(y_i)-F(y_{i-1})}{y_i-y_{i-1}}\Big)\int_0^{a_2}f(x_1^{i-1}(x_2),x_2)dx_2\Bigg)\vspace{3pt}\\
        \ge&- \frac{F(y_{i+1})-F(y_i)}{y_{i+1}-y_i}\|f_{x_1}\|_\infty\Big(\frac{F(y_{i+1})-F(y_i)}{y_{i+1}-y_i}-\frac{F(y_{i})-F(y_{i-1})}{y_{i}-y_{i-1}}\Big)\frac{a_2^2}{2}\vspace{3pt}\\
        &-\|f\|_\infty a_2^2\Big(\frac{F(y_{i+1})-F(y_i)}{y_{i+1}-y_i}
        -\frac{F(y_i)-F(y_{i-1})}{y_i-y_{i-1}}\Big),
        \end{array}\]
and we use $a_2^3\le a_2^2$ as $a_2\le1.$

 For the last term in \eqref{eqn: derivative computation}, we have
\[\hspace{-6pt}\begin{array}{ll}
    &\hspace{-8pt}\frac{c(z_{i+1})-c(z_i)}{y_{i+1}-y_i}\int_0^{a_2}f(x_1^{i}(x_2),x_2)dx_2-\frac{c(z_{i})-c(z_{i-1})}{y_i-y_{i-1}}\int_0^{a_2}f(x_1^{i-1}(x_2),x_2)dx_2\vspace{3pt}\\
    =&\hspace{-8pt}\frac{c(z_{i+1})-c(z_i)}{y_{i+1}-y_i}\int_0^{a_2}(f(x_1^{i}(x_2),x_2)-f(x_1^{i-1}(x_2),x_2))dx_2\vspace{3pt}\\
    &\hspace{-8pt}+\Big(\frac{c(z_{i+1})-c(z_i)}{y_{i+1}-y_i}-\frac{c(z_{i})-c(z_{i-1})}{y_i-y_{i-1}}\Big)\int_0^{a_2}f(x_1^{i-1}(x_2),x_2)dx_2\vspace{3pt}\\
    \ge&\hspace{-8pt}-\|f_{x_1}\|_\infty\frac{c(z_{i+1})-c(z_i)}{y_{i+1}-y_i}\Big(\frac{F(y_{i+1})-F(y_i)}{y_{i+1}-y_i}-\frac{F(y_{i})-F(y_{i-1})}{y_{i}-y_{i-1}}\Big)\frac{a_2^2}{2}\hspace{-3pt}+\hspace{-3pt}\alpha a_2\Big(\frac{c(z_{i+1})-c(z_i)}{y_{i+1}-y_i}-\frac{c(z_{i})-c(z_{i-1})}{y_i-y_{i-1}}\Big)
\end{array}\]

Using these bounds on (\ref{derivative}), we get
\[\hspace{-0.3cm}\begin{array}{ll}
&\frac{d}{d\epsilon}\Big|_{\epsilon=0} \mathcal P(v^\epsilon) \ge\vspace{3pt}\\
&\hspace{-0.3cm} -a_2^2\Big(\frac{F(y_{i+1})-F(y_i)}{y_{i+1}-y_i}-\frac{F(y_i)-F(y_{i-1})}{y_i-y_{i-1}}\Big)\Big(\frac{3}{2}\|f\|_\infty+\frac{\|f_{x_1}\|_\infty}{2}\Big(1+\frac{F(y_{i+1})-F(y_i)}{y_{i+1}-y_i}+\frac{c(z_{i+1})-c(z_i)}{y_{i+1}-y_i}\Big)\Big)\vspace{3pt}\\
&\hspace{-0.3cm}+\alpha a_2\Big(\frac{c(z_{i+1})-c(z_i)}{y_{i+1}-y_i}-\frac{c(z_{i})-c(z_{i-1})}{y_i-y_{i-1}}\Big)>0 \vspace{3pt}\\
\end{array}\]
by the fact that $a_2^2\le a_2$ and condition \ref{condition3}. This means decreasing the price of the $i$th good leads to a strictly larger profit.  Therefore, the original pricing schedule cannot be optimal.\end{proof}

\end{document}
